\uuid{kUlH}
\titre{Système linéaire/Méthode du Pivot de Gauss}
\auteur{hanani ; mzari}
\organisation{exo7}
\licence{CC BY-NC-SA 4.0}
\datecreate{2020-01-20}
\niveau{L1}
\module{Algèbre}
\chapitre{Déterminant, système linéaire}
\sousChapitre{Système linéaire, rang}
\difficulte{1}
\theme{Systèmes d'équations linéaires, Système linéaire, Méthode du Pivot de Gauss}
\exo7qcmid{lille-2.systemes.005}

\contenu{
\begin{question}
\qclassification{200.03}
\qtags{motcle=Système linéaire/Méthode du Pivot de Gauss}
On considère le système d'équations, d'inconnue $(x,y,z)\in \Rr^3$ :
$$(\mathtt{S}) \; \left\{\begin{array}{rcc}
x-y+z&=&1\\
2x-3y+4z&=&1\\
y+z&=&1.\\
\end{array}\right.$$
Quelles sont les assertions vraies ?
\begin{answers}
\good{$(\mathtt{S}) \Leftrightarrow  \left\{\begin{array}{rcc}
x-y+z&=&1\\
y-2z&=&1\\
z&=&0.\end{array}\right.$}
\good{Les équations de $(\mathtt{S})$ sont celles de trois plans}.
\good{$(\mathtt{S})$ admet une unique solution}.
\bad{$(\mathtt{S})$ n'admet pas de solution}.
\end{answers}
\begin{explanations}
$$(\mathtt{S}) \Leftrightarrow  \left\{\begin{array}{rcc}
x-y+z&=&1\\
y-2z&=&1\\
z&=&0.\\
\end{array}\right.$$
Donc $(\mathtt{S})$ admet une unique solution : $(2,1,0)$.
\end{explanations}
\end{question}
}
